\section{Câu 1. Mô hình hóa bài toán dưới dạng CSP: $P=(X,D,C)$}

\subsection{Tập biến $X$ và miền giá trị $D$}

Xét sáu nhân vật Alice, Bob, Carol, Dave, Eve và Frank ngồi quanh một bàn
tròn có sáu ghế. Gọi
\[
    X=\{A,B,C,D,E,F\}
\]
là tập biến, trong đó mỗi biến biểu diễn vị trí của một nhân vật:

\begin{center}
\begin{tabular}{c|cccccc}
    \toprule
    Biến & $A$ & $B$ & $C$ & $D$ & $E$ & $F$ \\
    \midrule
    Nhân vật & Alice & Bob & Carol & Dave & Eve & Frank \\
    \bottomrule
\end{tabular}
\end{center}

Các ghế được đánh số theo chiều kim đồng hồ từ 1 đến 6. Miền giá trị ban đầu
của mỗi biến là
\[
    D_A=D_B=D_C=D_D=D_E=D_F=\{1,2,3,4,5,6\}.
\]

Trên bàn tròn sáu ghế, hai ghế đối diện nhau khi hiệu tuyệt đối bằng 3; hai
ghế kề nhau khi hiệu tuyệt đối bằng 1 hoặc bằng 5.

\subsection{Tập ràng buộc $C$}

Mỗi ràng buộc được mô tả bởi cặp $(\text{scope},\text{relation})$. Các ràng
buộc của bài toán như sau.

\subsubsection{Ràng buộc một ngôi và các ràng buộc quan hệ}

\begin{enumerate}
    \item Alice ngồi ở ghế số 1:
    \[
        C_1=((A),\{1\}).
    \]

    \item Alice ngồi đối diện Bob. Quan hệ đối diện là
    \[
        R_{\mathrm{doi\ dien}}=
        \{(1,4),(2,5),(3,6),(4,1),(5,2),(6,3)\},
    \]
    nên
    \[
        C_2=((A,B),R_{\mathrm{doi\ dien}}).
    \]

    \item Frank không ngồi cạnh Alice. Quan hệ không kề được viết là
    \[
        R_{\mathrm{khong\ ke}}=
        \{(f,a)\in D_F\times D_A\mid |f-a|\notin\{1,5\}\},
    \]
    nên
    \[
        C_3=((F,A),R_{\mathrm{khong\ ke}}).
    \]

    \item Dave ngồi cạnh Eve:
    \[
        C_4=((D,E),R_{\mathrm{ke}}).
    \]

    \item Dave ngồi cạnh Bob:
    \[
        C_5=((D,B),R_{\mathrm{ke}}).
    \]
\end{enumerate}

Trong đó quan hệ hai ghế kề nhau là
\[
\begin{split}
R_{\mathrm{ke}}=\{&(1,2),(1,6),(2,1),(2,3),(3,2),(3,4),\\
                    &(4,3),(4,5),(5,4),(5,6),(6,5),(6,1)\}.
\end{split}
\]

\subsubsection{Ràng buộc AllDiff}

Mỗi nhân vật phải ngồi ở một ghế khác nhau. Gọi
\[
    R_{\ne}=\{(u,v)\in\{1,\ldots,6\}^2\mid u\ne v\}.
\]
Quan hệ này tạo ra 15 ràng buộc hai ngôi cho mọi cặp biến khác nhau:
\[
\begin{array}{lll}
 C_6=((A,B),R_{\ne}), & C_7=((A,C),R_{\ne}), & C_8=((A,D),R_{\ne}),\\
 C_9=((A,E),R_{\ne}), & C_{10}=((A,F),R_{\ne}), & C_{11}=((B,C),R_{\ne}),\\
 C_{12}=((B,D),R_{\ne}), & C_{13}=((B,E),R_{\ne}), & C_{14}=((B,F),R_{\ne}),\\
 C_{15}=((C,D),R_{\ne}), & C_{16}=((C,E),R_{\ne}), & C_{17}=((C,F),R_{\ne}),\\
 C_{18}=((D,E),R_{\ne}), & C_{19}=((D,F),R_{\ne}), & C_{20}=((E,F),R_{\ne}).
\end{array}
\]

Như vậy, mô hình CSP hoàn chỉnh là
\[
    P=(X,D,C),\qquad C=\{C_1,C_2,\ldots,C_{20}\}.
\]

